\subsection{What the parameterization enforces}
Provided the decoder remains within a path, $\hat d_i\geq0$ guarantees monotone progress in arc length, but does not impose constant speed. Mapping a flat sequence near zero through the exponential and clamp yields a dwell. A positive excursion increases distance covered during the same number of time steps. Fluctuations in consecutive increments determine longitudinal acceleration. The representation makes reversal difficult even when reversing would be physically sensible. It is therefore best suited to movement along a forward driving path.

The upper clip at five is in log-distance units. It corresponds to an increment of $\exp(5)-1$ metres in one quarter-second interval, much larger than ordinary movement in the setting. It is a numerical bound. Participant-specific speed limits and realistic accelerations must still be checked empirically. A hard geometric constraint cannot substitute for evidence about the distribution of sampled motion.

There is also a subtle distinction between the route branch drawn by the planner and the portion visited during the fixed horizon. Two 4,096-metre paths can diverge only after the first two minutes of motion. They are different plans yet may give indistinguishable observed coordinates during the sampled interval. Conversely, two candidates can follow one geometric path and have large time-aligned separation if one has stopped. Reporting both route counts and timed-coordinate separation avoids treating every distinct plan as a distinct behavioral outcome.
